quaternionic surface $A/L$, a normal field $L'/\mathbb Q$, and the final
primitive weight-two newform $f$, after the finite twist has been
absorbed, with
\[
 [L:\mathbb Q]\le576,\quad d'=[L':\mathbb Q]\le2304R,\quad
 N\le CR^2,\quad \operatorname{rad}(N)\mid R.
\]
Its two-dimensional compatible representation $r_\ell$ satisfies
\[
 V_\ell(A)|_{G_{L'}}\simeq
       (r_\ell\oplus r_\ell)|_{G_{L'}},\qquad
 (\det r_\ell)^2|_{G_{L'}}=\chi_{\rm cyc}^2|_{G_{L'}}.
\]
These are (LR3)--(LR5) and the determinant calculation in the existing
modularity argument. There is no bounded-degree hypothesis on the
coefficient field of $f$.

Let $\epsilon$ be the nebentypus of this final form. In the Tate-module
convention, with arithmetic Frobenius, $\det r_\ell=\epsilon\chi_{\rm cyc}$;
see \cite[\S12.5, (12.5.1)--(12.5.2) and Remark 12.5.1,
pp.~119--120]{diamond-im}. Equivalently, at a good prime its
characteristic polynomial is $X^2-a_pX+p\epsilon(p)$.
A consistent contragredient convention replaces the finite character by
its inverse, with the same kernel and order. The character $\epsilon^2$
is trivial on $G_{L'}$, so its image order divides $d'$. Consequently
\[
 t=\operatorname{ord}(\epsilon)\le2d'\le4608R.
\]
This bounds the character order, not the coefficient-field degree.

Put $H=\ker\epsilon\subset(\mathbb Z/N\mathbb Z)^*$ and
\[
 \Gamma_H=\left\{
 \begin{pmatrix}a&b\\c&d\end{pmatrix}\in\Gamma_0(N):
 d\bmod N\in H\right\}.
\]
Since $\epsilon(-1)=1$, this group contains $-I$.
Let $X_H=X_1(N)/H$ for the diamond action and $J_H=\operatorname{Jac}(X_H)$.
The curve, diamond action and quotient are defined over $\mathbb Q$.
The rational model here parametrizes a point of exact order $N$, up to
$H$, and is the model specified by \cite[\S2.1, (2.1.8)]{mascot-moduli};
it introduces no cyclotomic field-degree factor.

Write $q:J_1(N)\to A_f$ for the newform quotient. Its pulled-back
differentials are spanned by all Galois conjugates of $f$; see
\cite[Theorem~1(i),(iv), p.~526]{shimura-factors} and
\cite[\S3.1]{baker-et-al-finiteness}. Every $h\in H$ fixes those differentials,
since its eigenvalue on each conjugate is a conjugate of $\epsilon(h)=1$;
the diamond convention is \cite[\S3.2.1]{baker-et-al-finiteness}.
Homomorphisms of abelian varieties in characteristic zero are determined
by their differentials. Hence $q\circ\langle h\rangle=q$.
For
\[
 e_H=\frac1{|H|}\sum_{h\in H}\langle h\rangle
       \in\operatorname{End}_{\mathbb Q}^0(J_1(N))
\]
we therefore have $qe_H=q$.

Let $H_{\rm eff}$ be the image of $H$ in
$\operatorname{Aut}(X_1(N))$, and let $\pi:X_1(N)\to X_H$ be the quotient.
On Jacobians,
\[
 \pi^*\pi_*=\sum_{h\in H_{\rm eff}}h,\qquad
 \pi_*\pi^*=[|H_{\rm eff}|].
\]
The averages over $H$ and $H_{\rm eff}$ agree. In particular the trivial
action of $-1$ is not counted twice in the first identity.
The image of $e_H$ is thus $\mathbb Q$-isogenous to $J_H$, and its
restriction of $q$ is surjective. Poincare reducibility makes $A_f$
a $\mathbb Q$-isogeny factor of $J_H$.

The accepted realization of $r_\ell$ in $A_f$ puts it in $V_\ell(J_H)$.
Over $L'$, Faltings' Hom theorem and extension of scalars give a nonzero
homomorphism $A\to J_H$. Absolute simplicity and Poincare reducibility
give an $L'$-isogeny $A\times Z\to J_H$. In particular
$g=\dim J_H\ge2$. The Jacobian has its canonical principal polarization;
no bound on a polarization degree of $A$ or $Z$ is being assumed.

